\section{Выполнение работы}

\subsection{Описание задачи}

\subsubsection{Выбор датасета}

Для рассмотрения задачи регрессии был выбран датасет \href{https://archive.ics.uci.edu/dataset/320/student+performance}{Student Performance} с репозитория UCI Machine Learning. Он имеет следующие характеристики:
\begin{itemize}
    \item количество признаков --- 30;
    \item количество объектов --- 649;
    \item год публикации --- 2014.
\end{itemize}

В датасете представлены результаты студентов среднего образования двух португальских школ. Данные включают оценки студентов, демографические, социальные и школьные характеристики, и были собраны с использованием школьных отчётов и анкет.

Предоставлены два набора данных, касающиеся успеваемости по двум различным предметам: Математика и Португальский язык. 

\textbf{Важное замечание от авторов:} целевой признак G3 имеет сильную корреляцию с признаками G2 и G1. Это происходит потому, что G3 является итоговой оценкой за год (выданной в 3-м периоде), в то время как G1 и G2 соответствуют оценкам за 1-й и 2-й периоды.

\subsubsection{EDA}

Идеи:

groupby по medu и sex

merge / join по school

resample по age (добавим искусственно дату рождения -> возраст)

\dots

one hot encoding по reason

\subsubsection{Формулировка задачи}

Регрессировать Total Alcohol Consumption вместо G3

\subsection{Обучение моделей}

\dots

\subsection{Ответы на вопросы}

\dots

\subsection{Математическое обоснование MSE}

\dots